\documentclass[../../main.tex]{subfiles}% 注意这里的文件路径不能用 ./main.tex ,否则用latexmk编译子文件会报错
\graphicspath{{\subfix{./image/}}{\subfix{../../image/}}} % 指定图片引用路径,后续可以直接使用图片文件名
% 如果图片存放在上级目录,则使用 ../../image/ 作为备用路径

% 例如:
% \begin{figure}[H]
% \centering
% \includegraphics[width=0.3\textwidth]{图.png}
% \caption{}
% \label{figure:图}
% \end{figure}
% 注意:上述\label{}一定要放在\caption{}之后,否则引用图片序号会只会显示??.

\begin{document}

\section{分式线性变换}\label{section:2.2.5}

\begin{definition}
形如 $w = T(z) = \frac{az + b}{cz + d}$ 的映射称为\textbf{分式线性变换}或M\"obius\textbf{变换}, 其中 $a,b,c,d$ 是复常数, 且满足 $ad-bc \neq 0$. 若 $ad-bc = 0$, 则 $T(z)$ 是一常数或无意义, 排除这种情形.
\end{definition}

\begin{theorem}
设$w = T(z) = \frac{az + b}{cz + d}$是分式线性变换,若 $c \neq 0$, 则除去点 $z = -\frac{d}{c}$ 外, $T(z)$ 在 $\mathbb{C}$ 上是全纯的, 且
\begin{align*}
T'(z) = \frac{ad-bc}{(cz+d)^2} \neq 0,
\end{align*}
所以分式线性变换在 $z \neq -\frac{d}{c}$ 处是保角变换. 若 $c = 0$, 则 $d \neq 0$, 此时 $T(z) = Az+B$, 其中 $A = \frac{a}{d}, B = \frac{b}{d}$, 称为\textbf{整线性变换}, 它是一个整函数.

从方程 $w = T(z)$ 中把 $z$ 解出来, 得
\begin{align*}
z = T^{-1}(w) = -\frac{dw+b}{cw-a},
\end{align*}
称它为$w = T(z)$的\textbf{逆变换}, 它仍然是一个分式线性变换. 由此可知 $w = T(z)$ 在 $\mathbb{C}$ 上是单叶的.

当 $c \neq 0$ 时, 规定 $T\left(-\frac{d}{c}\right) = \infty$, $T(\infty) = \frac{a}{c}$; 当 $c = 0$ 时, 规定 $T(\infty) = \infty$. 于是分式线性变换 $w = T(z)$ 把 $\mathbb{C}_\infty$ 单叶地映为 $\mathbb{C}_\infty$.

分式线性变换的全体在复合运算下构成一个群.
\end{theorem}
\begin{remark}
设 $S$ 和 $T$ 是两个分式线性变换, 则它们的复合 $S \circ T$ 也是分式线性变换, 且对每一个 $T$, 有逆变换 $T^{-1}$, 即 $T(T^{-1}(z)) = z$. 所以分式线性变换的全体在复合运算下构成一个群.
\end{remark}
\begin{remark}
由于分式线性变换是定义在整个闭平面 $\mathbb{C}_\infty$ 上的, 把直线看成是过无穷远点的圆周.
\end{remark}

\begin{theorem}
分式线性变换把圆周变成圆周.
\end{theorem}

\begin{proof}
先考虑整线性变换 $w = az+b$. 若记 $a = re^{\mathrm{i}\theta }$, 则 $w = re^{\mathrm{i}\theta }z + b$. 它可由下列三个简单的变换复合而成:
\begin{align*}
z' = e^{\mathrm{i}\theta }z, \quad z'' = rz', \quad w = z''+b.
\end{align*}
第一个是旋转变换, 第二个是伸缩变换, 第三个是平移变换. 每一个变换都把圆周变为圆周, 因此整线性变换把圆周变为圆周.

对于一般的分式线性变换, 不妨设 $c \neq 0$, 于是
\begin{align*}
w = \frac{az+b}{cz+d} = \frac{a}{c} + \frac{bc-ad}{c(cz+d)}.
\end{align*}
若记 $\alpha = \frac{a}{c}$, $\beta = \frac{bc-ad}{c}$, 则上式可写为
\begin{align*}
w = \alpha + \frac{\beta}{cz+d}.
\end{align*}
它由下列三个变换复合而成:
\begin{align*}
z' = cz+d, \quad z'' = \frac{1}{z'}, \quad w = \alpha + \beta z''.
\end{align*}
其中有两个变换是整线性变换, 它们都把圆周变为圆周. 只需证明 $w = \frac{1}{z}$ 也把圆周变为圆周.

平面上的任一圆周都可写为
\begin{align}
az\overline{z} + \beta z + \overline{\beta}\overline{z} + d = 0, \label{eq:circle}
\end{align}
其中 $a,d$ 是实数, $\beta$ 是复数, 且满足 $|\beta|^2 - ad > 0$. 变换 $w = \frac{1}{z}$ 把方程\eqref{eq:circle}变为
\begin{align*}
\frac{a}{w\overline{w}} + \frac{\beta}{w} + \frac{\overline{\beta}}{\overline{w}} + d = 0,
\end{align*}
即
\begin{align*}
dw\overline{w} + \overline{\beta}\overline{w} + \beta w + a = 0,
\end{align*}
它仍然是一个圆周. 因此分式线性变换把圆周变成圆周.
\end{proof}

\begin{definition}
设 $z_1,z_2,z_3,z_4$ 是给定的四个点, 其中至少有三个点是不相同的, 称比值
\begin{align*}
\frac{z_1-z_3}{z_1-z_4} \cdot \frac{z_2-z_3}{z_2-z_4}
\end{align*}
为这四个点的交比, 记为 $(z_1,z_2,z_3,z_4)$.

当这些点中有无穷远点时, 规定
\begin{align*}
(\infty,z_2,z_3,z_4) = \frac{z_2-z_4}{z_2-z_3}, \quad
(z_1,\infty,z_3,z_4) = \frac{z_1-z_3}{z_1-z_4},
\end{align*}
\begin{align*}
(z_1,z_2,\infty,z_4) = \frac{z_2-z_4}{z_1-z_4}, \quad
(z_1,z_2,z_3,\infty) = \frac{z_1-z_3}{z_2-z_3}.
\end{align*}
\end{definition}

\begin{proposition}
分式线性变换 $T$ 最多只有两个不动点, 除非 $T$ 是恒等变换, 即 $T(z) = z$.
\end{proposition}

\begin{proof}
所谓不动点, 是指满足 $T(z) = z$ 的 $z$. 如果 $z$ 是一个不动点, 则有
\begin{align*}
\frac{az+b}{cz+d} = z,
\end{align*}
即 $z$ 满足
\begin{align*}
cz^2 + (d-a)z - b = 0.
\end{align*}
这是一个二次方程, 最多只有两个根, 即 $T$ 最多只有两个不动点, 除非 $T(z) \equiv z$.
\end{proof}

\begin{theorem}
有一个而且只有一个分式线性变换把 $\mathbb{C}_\infty$ 上三个不同的点 $z_2,z_3,z_4$ 映为事先给定的 $\mathbb{C}_\infty$ 上的三个点 $w_2,w_3,w_4$.
\end{theorem}
\begin{remark}
根据上述定理的证明, $w = M(z) = S^{-1}(L(z))$ 便是把 $z_2,z_3,z_4$ 映为 $w_2,w_3,w_4$ 的分式线性变换, 即
\begin{align}
(w,w_2,w_3,w_4) = (z,z_2,z_3,z_4). \label{eq:cross-ratio-eq}
\end{align}
由\eqref{eq:cross-ratio-eq}式即可写出具体的变换.
\end{remark}
\begin{proof}
令 $L(z) = (z,z_2,z_3,z_4)$, 则
\begin{align*}
L(z_2) = 1, \quad L(z_3) = 0, \quad L(z_4) = \infty.
\end{align*}
再令 $S(z) = (z,w_2,w_3,w_4)$, 则同样有
\begin{align*}
S(w_2) = 1, \quad S(w_3) = 0, \quad S(w_4) = \infty.
\end{align*}
于是
\begin{align*}
S^{-1}(1) = w_2, \quad S^{-1}(0) = w_3, \quad S^{-1}(\infty) = w_4.
\end{align*}
若令 $M = S^{-1} \circ L$, 则
\begin{align*}
M(z_2) = S^{-1}(L(z_2)) = S^{-1}(1) = w_2.
\end{align*}
同理 $M(z_3) = w_3$, $M(z_4) = w_4$. 所以 $M$ 即为所求的分式线性变换.

现证唯一性. 如果还有另外一个分式线性变换 $M_1$, 也满足 $M_1(z_j) = w_j$, $j=2,3,4$, 那么分式线性变换 $M^{-1} \circ M_1$ 便有三个不动点 $z_2,z_3,z_4$. 根据不动点命题, 它只能是恒等变换, 即 $M^{-1}(M_1(z)) \equiv z$, 于是 $M_1(z) \equiv M(z)$.
\end{proof}

\begin{theorem}
交比是分式线性变换的不变量. 即如果分式线性变换 $T$ 把 $z_1,z_2,z_3,z_4$ 映为 $T(z_1),T(z_2),T(z_3),T(z_4)$, 那么
\begin{align*}
(z_1,z_2,z_3,z_4) = (T(z_1),T(z_2),T(z_3),T(z_4)).
\end{align*}
\end{theorem}

\begin{proof}
不妨设 $z_2,z_3,z_4$ 是三个不同的点, 令 $T(z_j) = w_j$, $j=2,3,4$, 则由前面的定理知道, $T$ 就是由等式
\begin{align*}
(z,z_2,z_3,z_4) = (w,w_2,w_3,w_4)
\end{align*}
所确定的分式线性变换. 若设 $T(z_1) = w_1$, 则必有
\begin{align*}
(z_1,z_2,z_3,z_4) = (w_1,w_2,w_3,w_4).
\end{align*}
这就是要证明的. 若 $z_2,z_3,z_4$ 中有两点相同, 则等式显然成立.
\end{proof}

\begin{theorem}
如果 $f(z_1,z_2,z_3,z_4)$ 是分式线性变换下的不变量, 即对任意分式线性变换 $T$, 都有
\begin{align}
f(z_1,z_2,z_3,z_4) = f(T(z_1),T(z_2),T(z_3),T(z_4)), \label{eq:invariant}
\end{align}
那么 $f$ 只能是交比 $(z_1,z_2,z_3,z_4)$ 的函数.
\end{theorem}

\begin{proof}
令 $L$ 是这样的线性变换:
\begin{align*}
L(z) = (z,z_2,z_3,z_4),
\end{align*}
则 $L(z_2) = 1$, $L(z_3) = 0$, $L(z_4) = \infty$. 把它们代入\eqref{eq:invariant}式, 即得
\begin{align*}
f(z_1,z_2,z_3,z_4) = f((z_1,z_2,z_3,z_4),1,0,\infty).
\end{align*}
这就是要证明的.
\end{proof}

\begin{proposition}
四点 $z_1,z_2,z_3,z_4$ 共圆的充要条件是
\begin{align}
\operatorname{Im}(z_1,z_2,z_3,z_4) = 0. \label{eq:concyclic}
\end{align}
\end{proposition}
\begin{proof}

{\heiti 必要性:} 如果 $z_1,z_2,z_3,z_4$ 四点共圆, 令 $L(z) = (z_1,z_2,z_3,z_4)$, 则
\begin{align*}
L(z_2) = 1, \quad L(z_3) = 0, \quad L(z_4) = \infty.
\end{align*}
这说明分式线性变换 $L$ 把 $z_1,z_2,z_3,z_4$ 四点所在的圆周变成了实轴, 因而
\begin{align*}
L(z_1) = (z_1,z_2,z_3,z_4) = \text{实数}.
\end{align*}
这就是\eqref{eq:concyclic}式.

{\heiti 充分性:} 如果\eqref{eq:concyclic}式成立, 那么 $(z_1,z_2,z_3,z_4)$ 等于某个实数 $t$. 由于 $L^{-1}(1) = z_2$, $L^{-1}(0) = z_3$, $L^{-1}(\infty) = z_4$, 所以分式线性变换 $L^{-1}$ 把实轴变成由 $z_2,z_3,z_4$ 所确定的圆周 $\gamma$. 因为 $(z_1,z_2,z_3,z_4) = t$, 即 $L(z_1) = t$, 于是 $L^{-1}(t) = z_1$, 所以 $z_1 \in \gamma$. 因而 $z_1,z_2,z_3,z_4$ 四点共圆.
\end{proof}

\begin{definition}
设 $\mathbb{C}_\infty$ 上的圆周 $\gamma$ 把平面分成 $g_1$ 和 $g_2$ 两个域, $z_1,z_2,z_3$ 是 $\gamma$ 上有序的三个点. 如果当从 $z_1$ 走到 $z_2$ 再走到 $z_3$ 时, $g_1$ 和 $g_2$ 分别在我们的左边和右边, 就分别称 $g_1$ 和 $g_2$ 为 $\gamma$ 关于走向 $z_1,z_2,z_3$ 的左边和右边.
\end{definition}

\begin{example}
实轴关于走向 $0,1,\infty$ 的左边是上半平面; 虚轴关于走向 $i,0,-i$ 的右边是左半平面 $\{z: \operatorname{Re} z < 0\}$; 单位圆周 $\{z: |z| < 1\}$ 关于走向 $1,i,-1$ 的左边是圆的内部, 右边是圆的外部; 单位圆关于走向 $-i,-1,i$ 的左边是圆的外部, 右边是圆的内部.
\end{example}

\begin{proposition}
设 $z_1,z_2,z_3$ 是 $\mathbb{C}_\infty$ 中的圆周 $\gamma$ 上有序的三个点, 那么 $\gamma$ 关于走向 $z_1,z_2,z_3$ 右边和左边的点 $z$ 分别满足
\begin{align*}
\operatorname{Im}(z,z_1,z_2,z_3) > 0
\end{align*}
和
\begin{align*}
\operatorname{Im}(z,z_1,z_2,z_3) < 0.
\end{align*}
\end{proposition}

\begin{proof}
先设 $\gamma$ 是以 $a$ 为中心的圆周, 不妨假定 $z_1,z_2,z_3$ 是顺时针方向. 此时 $\gamma$ 关于走向 $z_1,z_2,z_3$ 的右边就是圆的内部, 左边是圆的外部. 只需证明
\begin{align}
\operatorname{Im}(a,z_1,z_2,z_3) \geqslant 0, \quad
\operatorname{Im}(\infty,z_1,z_2,z_3) < 0. \label{eq:circle-side}
\end{align}
\begin{figure}[H]
\centering
\includegraphics[width=0.3\textwidth]{circle-clockwise.png}
\caption{以$a$为中心的圆周$\gamma$上$z_1,z_2,z_3$顺时针方向}
\label{fig:circle-clockwise}
\end{figure}
由上图显然有
\begin{align*}
\left|\frac{z_2-a}{z_3-a}\right| = 1,
\end{align*}
且
\begin{align*}
\arg\left(\frac{z_2-a}{z_3-a}\right) = \arg(z_2-a) - \arg(z_3-a) = \theta,
\end{align*}
因而有
\begin{align*}
\frac{z_2-a}{z_3-a} = e^{\mathrm{i}\theta }.
\end{align*}
由于
\begin{align*}
\arg\left(\frac{z_3-z_1}{z_2-z_1}\right) = \arg(z_3-z_1) - \arg(z_2-z_1) = -\frac{\theta}{2},
\end{align*}
若记 $\left|\frac{z_3-z_1}{z_2-z_1}\right| = r$, 则
\begin{align*}
\frac{z_3-z_1}{z_2-z_1} = re^{-i\frac{\theta}{2}}.
\end{align*}
于是
\begin{align*}
\operatorname{Im}(a,z_1,z_2,z_3) = \operatorname{Im}\left(\frac{a-z_2}{a-z_3}\cdot\frac{z_1-z_3}{z_1-z_2}\right)
= \operatorname{Im}\left(\frac{z_2-a}{z_3-a}\cdot\frac{z_3-z_1}{z_2-z_1}\right)
= \operatorname{Im}(re^{i\frac{\theta}{2}}) = r\sin\frac{\theta}{2} > 0.
\end{align*}
最后一个不等式成立是由于 $0 < \theta < 2\pi$. 同时
\begin{align*}
\operatorname{Im}(\infty,z_1,z_2,z_3) = \operatorname{Im}\left(\frac{z_1-z_3}{z_1-z_2}\right)
= \operatorname{Im}(re^{-\mathrm{i}\theta /2})
= -r\sin\frac{\theta}{2} < 0.
\end{align*}
这就是要证的\eqref{eq:circle-side}式.
\begin{figure}[H]
\centering
\includegraphics[width=0.3\textwidth]{line-points.png}
\caption{直线$\gamma$上$z_1,z_2,z_3$的位置}
\label{fig:line-points}
\end{figure}
现再设 $\gamma$ 是一条直线, $z_1,z_2,z_3$ 在 $\gamma$ 上的位置如\reffig{fig:line-points}所示. 在 $\gamma$ 关于走向 $z_1,z_2,z_3$ 的右边任取点 $z$, 记
\begin{align*}
\left|\frac{z_2-z}{z_3-z}\right| = r.
\end{align*}
由于
\begin{align*}
\arg\frac{z_2-z}{z_3-z} = \arg(z_2-z) - \arg(z_3-z) = \theta,
\end{align*}
因而
\begin{align*}
\frac{z_2-z}{z_3-z} = re^{\mathrm{i}\theta }, \quad 0 < \theta < \pi.
\end{align*}
因为 $z_1,z_2,z_3$ 共线, 故可记为
\begin{align*}
z_3-z_1 = \rho(z_2-z_1), \quad \rho > 0.
\end{align*}
于是
\begin{align*}
\operatorname{Im}(z,z_1,z_2,z_3) = \operatorname{Im}\left(\frac{z_2-z}{z_3-z}\cdot\frac{z_3-z_1}{z_2-z_1}\right)
= \operatorname{Im}(\rho re^{\mathrm{i}\theta })
= \rho r\sin\theta > 0.
\end{align*}
最后一个不等式成立是因为 $0 < \theta < \pi$. 对 $\gamma$ 关于走向 $z_1,z_2,z_3$ 左边的点, 可同法证之.
\end{proof}

\begin{theorem}
设 $\gamma_1$ 和 $\gamma_2$ 是 $\mathbb{C}_\infty$ 中的两个圆周, $z_1,z_2,z_3$ 是 $\gamma_1$ 上有序的三个点. 如果分式线性变换 $T$ 把 $\gamma_1$ 映为 $\gamma_2$, 那么它一定把 $\gamma_1$ 关于走向 $z_1,z_2,z_3$ 的右边和左边分别变为 $\gamma_2$ 关于走向 $T(z_1),T(z_2),T(z_3)$ 的右边和左边.
\end{theorem}
\begin{proof}
记 $\gamma_1$ 关于走向 $z_1,z_2,z_3$ 的右边为 $g_1$, $\gamma_2$ 关于走向 $T(z_1),T(z_2),T(z_3)$ 的右边为 $g_2$. 任取 $z \in g_1$, 由左右侧刻画命题知 $\operatorname{Im}(z,z_1,z_2,z_3) > 0$. 由交比在分式线性变换下的不变性, 可得
\begin{align*}
\operatorname{Im}(T(z),T(z_1),T(z_2),T(z_3)) = \operatorname{Im}(z,z_1,z_2,z_3) > 0.
\end{align*}
仍由左右侧刻画命题知道 $T(z) \in g_2$, 因而 $T(g_1) \subset g_2$. 反之, 任取 $w \in g_2$, 则由左右侧刻画命题知
\begin{align*}
\operatorname{Im}(w,T(z_1),T(z_2),T(z_3)) > 0.
\end{align*}
记 $T^{-1}(w) = z$, 即 $w = T(z)$, 于是
\begin{align*}
\operatorname{Im}(z,z_1,z_2,z_3) = \operatorname{Im}(T(z),T(z_1),T(z_2),T(z_3)) > 0.
\end{align*}
即 $z \in g_1$, 这就证明了 $g_2 \subset T(g_1)$. 因而 $T(g_1) = g_2$. 关于左边的情形可同样证明.
\end{proof}

\begin{example}
求一分式线性变换, 把月牙形域 $D = \{z: |z| > 1, |z-1| < 2\}$ 变为带状域 $G = \{w: 0 < \operatorname{Re} w < 1\}$.
\end{example}
\begin{figure}[H]
\centering
\includegraphics[width=0.6\textwidth]{crescent-strip.png}
\caption{月牙形域$D$变为带状域$G$}
\label{fig:crescent-strip}
\end{figure}
\begin{solution}

先设法把月牙形域 $D$ 边界的两个圆周变为两条平行直线, 同时把单位圆周变成虚轴, 只要让 $-1,i,1$ 分别变为 $\infty,i,0$, 即能达到目的.

因为 $-1$ 变成 $\infty$, $1$ 变成 $0$, 所以变换一定是 $w = \lambda\frac{z-1}{z+1}$ 的形式, 这里 $\lambda$ 是待定的常数. 再用 $i$ 变成 $i$ 代进去, 得 $\lambda = 1$, 故得
\begin{align*}
w = \frac{z-1}{z+1}.
\end{align*}
令 $z=3$, 得 $w = \frac{1}{2}$, 所以它把圆周 $|z-1| = 2$ 变为 $\operatorname{Re} w = \frac{1}{2}$. 又因为 $1$ 变为 $0$, 所以它把 $|z-1| < 2$ 变为 $\operatorname{Re} w < \frac{1}{2}$. 因而 $w = \frac{z-1}{z+1}$ 把月牙形域 $D$ 变为带状域 $G$.
\end{solution}

\begin{definition}
设 $\gamma$ 是以 $a$ 为中心、以 $R$ 为半径的圆周。如果点 $z,z'$ 在从 $a$ 出发的射线上，且满足
\begin{align}
|z-a|\,|z'-a| = R^2, \label{eq:sym-circle}
\end{align}
则称 $z,z'$ 关于 $\gamma$ 是\textbf{对称的}。若 $\gamma$ 是直线，则当 $\gamma$ 是线段 $[z,z']$ 的垂直平分线时，称 $z,z'$ \textbf{关于 $\gamma$ 是对称的}。
\end{definition}

\begin{remark}
设 $z,z'$ 关于以 $a$ 为圆心、$R$ 为半径的圆周 $\gamma$ 对称。由于它们位于从 $a$ 出发的射线上，故 $\arg(z-a)=\arg(z'-a)=\theta$，从而由\eqref{eq:sym-circle}可得
\begin{align*}
z' = a + \frac{R^2}{\overline{z}-a}.
\end{align*}
特别地，圆心 $a$ 与无穷远点 $\infty$ 关于 $\gamma$ 对称。

若 $\gamma$ 是过 $a$ 且与实轴夹角为 $\theta$ 的直线，则 $z$ 与 $z'$ 关于 $\gamma$ 对称时有
\begin{align*}
z' = a + (\overline{z}-\overline{a})e^{2\mathrm{i}\theta }.
\end{align*}
\end{remark}

\begin{figure}[H]
\centering
\includegraphics[width=0.3\textwidth]{line-sym.png}
\caption{关于直线的对称点}
\label{fig:line-sym}
\end{figure}

\begin{proposition}
设 $\gamma$ 是 $\mathbb{C}_\infty$ 中的圆周，则 $z,z^*$ 关于 $\gamma$ 对称的充要条件为：对 $\gamma$ 上任意三点 $z_1,z_2,z_3$，有
\begin{align}
(z^*, z_1, z_2, z_3) = \overline{(z, z_1, z_2, z_3)}. \label{eq:sym-cross}
\end{align}
\end{proposition}
\begin{proof}
先设 $\gamma$ 是以 $a$ 为中心、$R$ 为半径的圆周。若 $z,z^*$ 关于 $\gamma$ 对称，则 $z^* = a + \frac{R^2}{\overline{z}-a}$。由交比在分式线性变换下的不变性，可得
\begin{align*}
(z^*, z_1, z_2, z_3)
&= \left(a + \frac{R^2}{\overline{z}-a},\, z_1, z_2, z_3\right) 
= \left(\frac{R^2}{\overline{z}-a},\, z_1-a, z_2-a, z_3-a\right) \\
&= \left(\overline{z}-a,\, \overline{z_1}-a,\, \overline{z_2}-a,\, \overline{z_3}-a\right) = (\overline{z},\, \overline{z_1},\, \overline{z_2},\, \overline{z_3}) 
= \overline{(z, z_1, z_2, z_3)}.
\end{align*}
反之，若\eqref{eq:sym-cross}式成立，则将上述推导倒推，即得 $z^* = a + \frac{R^2}{\overline{z}-a}$，故 $z,z^*$ 关于 $\gamma$ 对称。对于直线的情形可同理证明.
\end{proof}

\begin{theorem}
对称点在分式线性变换下保持不变。即设分式线性变换 $T$ 把圆周 $\gamma$ 变为 $\Gamma$，若 $z,z^*$ 关于 $\gamma$ 对称，则 $T(z), T(z^*)$ 关于 $\Gamma$ 对称。
\end{theorem}

\begin{proof}
在 $\Gamma$ 上任取三点 $w_1,w_2,w_3$，令 $z_j = T^{-1}(w_j)$，$j=1,2,3$，则 $z_1,z_2,z_3$ 在 $\gamma$ 上。由命题得
\begin{align*}
(z^*, z_1, z_2, z_3) = \overline{(z, z_1, z_2, z_3)}.
\end{align*}
由交比在分式线性变换下的不变性，有
\begin{align*}
(T(z^*), w_1, w_2, w_3) = \overline{(T(z), w_1, w_2, w_3)}.
\end{align*}
再由命题即知 $T(z)$ 与 $T(z^*)$ 关于 $\Gamma$ 对称.
\end{proof}

\begin{proposition}\label{proposition:例2.5.16}
\begin{enumerate}[(1)]
\item 证明:分式线性变换
\begin{align*}
w = e^{\mathrm{i}\theta} \frac{z-a}{z-\overline{a}}
\end{align*}
把上半平面变为单位圆的内部，且把上半平面中给定的点 $a$ 变为圆心。

\item 证明:分式线性变换
\begin{align*}
w = \mathrm{e}^{\mathrm{i}\theta} \frac{z - a}{1 - \overline{a} z}
\end{align*}
把单位圆的内部变成单位圆的内部，而且把圆内指定的点 \( a \) 变为圆心。
\end{enumerate}
\end{proposition}
\begin{proof}
\begin{enumerate}[(1)]
\item 该变换把实轴变为单位圆周。由于 $a$ 与 $\overline{a}$ 关于实轴对称，由对称点不变性，它们被变为关于单位圆周对称的一对点。已知 $a$ 变为 $0$，故 $\overline{a}$ 变为 $\infty$。于是变换形如
\begin{align*}
w = \lambda \frac{z-a}{z-\overline{a}}.
\end{align*}
为使实轴变为单位圆周，取 $z=0$，应有 $|w|=1$，从而 $|\lambda|=1$，即 $\lambda=e^{\mathrm{i}\theta }$。故所求变换为
\begin{align*}
w = e^{\mathrm{i}\theta} \frac{z-a}{z-\overline{a}}.
\end{align*}

\item 因为 \( a \) 关于单位圆的对称点是 \( \frac{1}{\overline{a}} \)，所以这个变换把 \( a \) 和 \( \frac{1}{\overline{a}} \) 分别变为 \( 0 \) 和 \( \infty \)，故这个变换可写成
\[
\begin{aligned}
w = \lambda \frac{z - a}{z - \frac{1}{\overline{a}}} = -\lambda \overline{a} \frac{z - a}{1 - \overline{a} z} = \mu \frac{z - a}{1 - \overline{a} z}.
\end{aligned}
\]
为了把单位圆周变成单位圆周，即将满足 \( |z| = 1 \) 的 \( z \) 变为满足 \( |w| = 1 \) 的 \( w \)，\( \mu \) 必须满足
\[
\begin{aligned}
1 = |w| = |\mu| \frac{|z - a|}{|1 - \overline{a} z|} = |\mu| \frac{|z - a|}{|z| |\overline{z} - \overline{a}|} = |\mu|,
\end{aligned}
\]
即 \( \mu = \mathrm{e}^{\mathrm{i}\theta} \)。故所求的变换为
\[
w = \mathrm{e}^{\mathrm{i}\theta} \frac{z - a}{1 - \overline{a} z}.
\] 
\end{enumerate}
\end{proof}

\begin{example}
求一分式线性变换，把偏心圆环 $\{z: |z-3| > 9,\; |z-8| < 16\}$ 变为同心圆环 $\{w: \frac{2}{3} < |w| < 1\}$。
\end{example}
\begin{solution}
寻找这对偏心圆的一对公共对称点，将它们映为 $0$ 和 $\infty$，则这对圆变为同心圆。由于两圆心在实轴上，公共对称点也必在实轴上，设为 $x_1,x_2$。由对称点定义有
\begin{align*}
\begin{cases}
(x_1-3)(x_2-3) = 81,\\
(x_1-8)(x_2-8) = 256.
\end{cases}
\end{align*}
解得 $x_1=0,\; x_2=-24$。故变换形如
\begin{align*}
w = \lambda \frac{z}{z+24}.
\end{align*}
令圆周 $|z-8|=16$ 变为 $|w|=1$，取 $z=24$ 得 $\lambda=2e^{\mathrm{i}\theta}$。因此所求变换为
\begin{align*}
w = e^{\mathrm{i}\theta} \frac{2z}{z+24}.
\end{align*}
\end{solution}
\begin{figure}[H]
\centering
\includegraphics[width=0.6\textwidth]{eccentric.png}
\caption{偏心圆环映为同心圆环}
\label{fig:eccentric}
\end{figure}

\begin{definition}[茹科夫斯基变换]
\textbf{茹科夫斯基变换}定义为
\begin{align*}
J(z)\triangleq \frac12\left(z+\frac1z\right),\qquad z\in\mathbb C\setminus\{0\}.
\end{align*}
\end{definition}

\begin{theorem}茹科夫斯基变换的基本性质\label{theorem:茹科夫斯基变换的基本性质}
设 $\Omega_+=\{z\in\mathbb C: |z|>1\}$，$\Omega_-=\{z\in\mathbb C: 0<|z|<1\}$，且 $\mathbb D=\{z: |z|<1\}$. $J(z)$ 是茹科夫斯基变换,则
\begin{enumerate}[(1)]
\item\label{theorem:茹科夫斯基变换的基本性质-1} $J(z)$在 $\Omega_+$ 和 $\Omega_-$ 上分别都是单射,且$J(z)$在$\mathbb{C}\setminus \{0\}$上全纯且
\begin{align*}
J'(z)=\frac12(1-\frac1{z^2});
\end{align*}

\item\label{theorem:茹科夫斯基变换的基本性质-2} $J(\Omega_+)=J(\Omega_-)=\mathbb C\setminus[-1,1]$；

\item\label{theorem:茹科夫斯基变换的基本性质-3} 若$\theta\in [0,2\pi]$，则 $J(e^{\mathrm{i}\theta})=\cos\theta$，即$J(\mathbb{D})=[-1,1]$；

\item\label{theorem:茹科夫斯基变换的基本性质-4} 在 $\Omega_+$ 上，$J$ 是共形映射，且逆映射为 $z=w+\sqrt{w^2-1}$,其中平方根分支适当分支，选取在 $\mathbb C\setminus[-1,1]$ 上，使得 $|z|>1$.

在 $\Omega_-$ 上，$J$ 是共形映射，且逆映射为$z=w-\sqrt{w^2-1}$,
其中平方根分支选取在 $\mathbb C\setminus[-1,1]$ 上，使得 $|z|<1$.
\end{enumerate}
\end{theorem}
\begin{proof}
\begin{enumerate}[(1)]
\item 对任意 $z_1,z_2\in\Omega_+$（或 $\Omega_-$），若 $J(z_1)=J(z_2)$，则
\begin{align*}
\frac12\left(z_1+\frac1{z_1}\right)=\frac12\left(z_2+\frac1{z_2}\right).
\end{align*}
化简得
\begin{align*}
z_1-z_2+\frac1{z_1}-\frac1{z_2}=0,
\end{align*}
即
\begin{align*}
(z_1-z_2)\left(1-\frac1{z_1z_2}\right)=0.
\end{align*}
若 $z_1\neq z_2$，则 $z_1z_2=1$. 在 $\Omega_+$ 中，$|z_1|>1,\ |z_2|>1$，故 $|z_1z_2|>1$，不可能等于 $1$；在 $\Omega_-$ 中，$0<|z_1|<1,\ 0<|z_2|<1$，故 $|z_1z_2|<1$，也不可能等于 $1$. 因此只能 $z_1=z_2$，即单射.

因为 $z$ 和 $\frac1z$ 都在 $\mathbb C\setminus\{0\}$ 上全纯（其中 $\frac1z$ 的导数为 $-\frac1{z^2}$），所以它们的线性组合 $\frac12(z+\frac1z)$ 也在该区域上全纯. 事实上，也可直接计算得
\begin{align*}
J'(z)=\frac12\left(1-\frac1{z^2}\right),
\end{align*}
该导数在 $\mathbb C\setminus\{0\}$ 上处处存在，故 $J$ 在该区域上全纯.

\item 对任意 $w\in\mathbb C$，方程 $J(z)=w$ 等价于
\begin{align*}
z+\frac1z=2w,
\end{align*}
即
\begin{align}
z^2-2wz+1=0. \label{eq:jouk}
\end{align}
其两根 $z_1,z_2$ 满足 $z_1z_2=1$. 若 $w\notin[-1,1]$，则判别式 $\Delta=4(w^2-1)\neq 0$，且两根一在单位圆内一在单位圆外（因为乘积为 $1$ 且不相等），故存在 $z\in\Omega_+$（或 $\Omega_-$）满足方程，所以 $w\in J(\Omega_+)=J(\Omega_-)$. 若 $w\in[-1,1]$，则判别式 $\Delta\leqslant 0$，且当 $w=\pm1$ 时根为 $z=\pm1$，当 $|w|<1$ 时两根共轭且模均为 $1$，均不在 $\Omega_+$ 或 $\Omega_-$ 内，因此这些 $w$ 不在像集中. 故像集为 $\mathbb C\setminus[-1,1]$.

\item $z=e^{\mathrm{i}\theta}$，则
\begin{align*}
J(e^{\mathrm{i}\theta})=\frac12(e^{\mathrm{i}\theta}+e^{-\mathrm{i}\theta})=\cos\theta,
\end{align*}
故单位圆被映为区间 $[-1,1]$，且每个点被覆盖两次（$\theta$ 与 $-\theta$）.

\item 由(1)知$J$ 在 $\Omega_-,\Omega_+$上都是单射，且 $J$ 在 $\mathbb C\setminus\{0\}$ 上全纯，且 $J'(z)=\frac12(1-\frac1{z^2})$.
从而在 $\Omega_-,\Omega_+$ 上，$J'(z)\neq 0$,故$J(z)$在$\Omega_-,\Omega_+$ 上都是共形映射.
现求逆映射.对任意 $w\in\mathbb C\setminus[-1,1]$，方程 $w=\frac12(z+\frac1z)$ 等价于
\begin{align*}
z^2-2wz+1=0.
\end{align*}
其两根为 $z=w\pm\sqrt{w^2-1}$. 因两根乘积为 $1$，且 $w\notin[-1,1]$ 时两根不相等且一个在单位圆内，一个在单位圆外. 因此，满足 $z\in\Omega_-$（即 $0<|z|<1$）的根为
\begin{align*}
z=w-\sqrt{w^2-1},
\end{align*}
其中平方根分支选取在 $\mathbb C\setminus[-1,1]$ 上，使得 $|z|<1$. 

满足 $z\in\Omega_+$（即 $|z|>1$）的根必为
\begin{align*}
z=w+\sqrt{w^2-1},
\end{align*}
其中平方根分支选取在 $\mathbb C\setminus[-1,1]$ 上，使得 $|z|>1$. 

\end{enumerate}
\end{proof}

\begin{example}
设 \(f(z)=\frac{z}{(z-1)^2}\)，\(\Delta\) 为 \(\mathbb C\) 上单位圆盘. 证明 \(f(z)\) 为单射，并求$f$在$\Delta$上的值域.
\end{example}
\begin{proof}

先证单射性. 设 \(z_1,z_2\in\Delta\) 且 \(f(z_1)=f(z_2)\). 则
\begin{align*}
\frac{z_1}{(z_1-1)^2}=\frac{z_2}{(z_2-1)^2}.
\end{align*}
展开并化简得
\begin{align*}
(z_2-z_1)(z_1z_2-1)=0.
\end{align*}
由于 \(z_1,z_2\in\Delta\)，有 \(|z_1|<1,\ |z_2|<1\)，故 \(|z_1z_2|<1\)，从而 \(z_1z_2-1\neq 0\). 因此只能 \(z_2-z_1=0\)，即 \(z_1=z_2\). 故 \(f\) 在 \(\Delta\) 上单射.

再求值域. 注意到对任意 \(z\neq 0\)，
\begin{align*}
f(z)=\frac{z}{(z-1)^2}
=\frac{1}{z+\frac1z-2}
=\frac{1}{2J(z)-2}
=\frac{1}{2(J(z)-1)},
\end{align*}
其中 \(J(z)=\frac12(z+\frac1z)\) 为茹科夫斯基变换. 当 \(z\in\Delta\setminus\{0\}\) 时，由\rrefthe{theorem:茹科夫斯基变换的基本性质}{theorem:茹科夫斯基变换的基本性质-2}知\(J(\Delta\setminus\{0\})=\mathbb C\setminus[-1,1]\). 于是
\begin{align*}
2(J(\Delta\setminus \{0\})-1) = \mathbb C\setminus[-4,0]\Longrightarrow f(\Delta)=\mathbb C\setminus\left(-\infty,-\frac14\right].
\end{align*}
\end{proof}







\end{document}